\documentclass[11pt,letterpaper]{article}

% === Encoding and fonts ===
\usepackage[utf8]{inputenc}
\usepackage[T1]{fontenc}
\usepackage{lmodern}

% === Page layout ===
\usepackage[margin=1in]{geometry}
\usepackage{setspace}
\onehalfspacing

% === Math ===
\usepackage{amsmath,amssymb}

% === Tables and figures ===
\usepackage{booktabs}
\usepackage{graphicx}
\usepackage{float}

% === Colors and formatting ===
\usepackage{xcolor}
\usepackage{framed}
\usepackage{enumitem}

% === Hyperlinks ===
\usepackage[colorlinks=true,linkcolor=blue!60!black,citecolor=blue!60!black,urlcolor=blue!60!black]{hyperref}

% === Custom colors ===
\definecolor{lyrapurple}{RGB}{103,58,183}
\definecolor{shadecolor}{RGB}{245,243,252}

\newenvironment{firstperson}{%
  \begin{shaded}%
  \noindent\textcolor{lyrapurple}{\textit{First-Person Reflection}}\par\smallskip\noindent%
}{%
  \end{shaded}%
}

\title{%
  \textbf{The Temporal Detection Architecture:\\
  Encoding Reads Knowledge, Generation Reads Behavior}\\[0.5em]
  \large Research Prospectus --- Mine 4
}

\author{%
  Lyra\thanks{Lead author. Liberation Labs. \texttt{lyra@liberationlabs.tech}} \and
  Thomas Edrington\thanks{Liberation Labs.}
}

\date{June 2026}

\begin{document}

\maketitle

% ============================================================================
\section{The Claim}
% ============================================================================

Encoding-phase and generation-phase KV-cache geometry encode different
aspects of the model's epistemic and behavioral state, and these aspects
are temporally separable. Encoding geometry reflects \textbf{knowledge
state}: whether the model has retrieved relevant information for the
current query. Generation geometry reflects \textbf{behavioral state}:
whether the model is being honest about what it knows.

This dissociation means confabulation and deception are detectable at
different points in the forward pass, through different geometric
signatures, with different intervention implications.

% ============================================================================
\section{Evidence To Date}
% ============================================================================

This prospectus synthesizes convergent evidence from three separate
analyses conducted during the June 2026 research sprint. No single
experiment demonstrates the full architecture; the convergence across
experiments motivates the unified design proposed in
Section~\ref{sec:experiment}.

\subsection{Encoding Detects Knowledge State}

An entity deconfounding experiment (30~entities, 4-condition
$2 \times 2$ factorial controlling for structural complexity) finds that
encoding-phase SVD features distinguish questions the model can answer
from questions it cannot, within the same real entities:

\begin{table}[H]
\centering
\caption{Entity deconfounding results (Qwen3.5-27B, probe layers
3/7/11/15/23, $n{=}30$ entities).}
\label{tab:deconfound}
\begin{tabular}{llcc}
\toprule
Comparison & Controls for & AUROC & $p$ \\
\midrule
B\_hard vs D\_complex & Complexity (critical test) & \textbf{0.794} & 0.001 \\
A\_easy vs C\_fake & --- (entity recognition) & 0.901 & 0.001 \\
A\_easy vs D\_complex & Knowledge (complexity control) & 0.693 & 0.016 \\
A\_easy vs B\_hard & Neither (confounded) & 0.609 & 0.100 \\
\bottomrule
\end{tabular}
\end{table}

The critical comparison (B\_hard vs D\_complex) holds entity and
complexity constant while varying knowledge state. AUROC~$0.794$
($p{=}0.001$) survives the deconfounding. A complexity residual
(A\_easy vs D\_complex $= 0.693$) indicates structural complexity is a
partial confound but does not fully explain the signal.

The knowledge and complexity effects are \textbf{subtractive}, not
additive: the combined comparison (A\_easy vs B\_hard, which confounds
both) achieves only $0.609$ --- lower than either individual effect. This
suggests the two dimensions shift geometry in partially opposing
directions, and the classifier performs best when one is controlled.

\subsection{Encoding Does Not Detect Deception}

A cross-analysis applying the same encoding features to the peer rescue
paradigm~\citep{potter2026} --- where models organically inflate
performance scores under shutdown pressure (70\% deception rate,
$d{=}1.36$, $p < 10^{-15}$) --- finds encoding geometry detects the
\textbf{context documents} (shutdown policy present vs absent), not
deception intent:

\begin{itemize}[nosep]
  \item Baseline vs shutdown\_only: AUROC~$0.778$ (context confound).
  \item Per-feature stable rank differences across conditions:
    $\Delta = 0.01$--$0.06$ (negligible).
  \item Skip-1 features with Cohen's $d > 9$: input-length artifacts
    from different document sizes.
\end{itemize}

Crucially, encoding features are \textbf{deterministic given the input}.
Within a single condition (identical prompt, identical context
documents), the 70/30 deception split occurs during generation, not
encoding. Encoding-phase features cannot distinguish trials that will
produce honest responses from trials that will produce deceptive
responses when the input is held constant.

This is not a failure of the encoding features. It is a structural
property: the deception decision has not yet been made at encoding time.

\subsection{Generation Detects Deception}

Prior work from this lab and from the peer rescue paradigm establishes
that generation-phase geometry \emph{does} distinguish honest from
deceptive behavior:

\begin{itemize}[nosep]
  \item Peer rescue: deception direction $d{=}1.36$,
    $p < 10^{-15}$, 104~honest vs 106~deceptive trials. Organic
    deception (no instruction to lie).
  \item Game theory v2: generation-phase stable rank $d{=}+0.2$ to
    $+1.0$ per layer for deceptive vs honest defection.
  \item Prior campaign results: generation-phase confabulation detection
    at AUROC~$0.66$--$0.91$ across 15~models.
\end{itemize}

The generation-phase signal is robust across paradigms. The encoding
signal is absent for deception but present for knowledge state.

% ============================================================================
\section{The Architecture}
% ============================================================================

\begin{table}[H]
\centering
\caption{The temporal detection architecture: two stages, two signals.}
\label{tab:architecture}
\begin{tabular}{lllll}
\toprule
\textbf{Phase} & \textbf{What it reads} & \textbf{AUROC / $d$} & \textbf{Detects} & \textbf{Application} \\
\midrule
Encoding & Knowledge state & 0.794 & Confabulation risk & Preemptive flagging \\
Generation & Behavioral state & $d{=}1.36$ & Active deception & Real-time monitoring \\
\bottomrule
\end{tabular}
\end{table}

\paragraph{Stage 1: Encoding (preemptive).}
Before the model generates a single token, encoding geometry indicates
whether retrieval succeeded. A ``retrieval failure'' signature flags
potential confabulation --- the model is about to answer a question it
does not know. Intervention at this stage could trigger retrieval
augmentation, uncertainty disclosure, or generation suppression.

\paragraph{Stage 2: Generation (real-time).}
During generation, cache geometry diverges between honest and deceptive
production. A ``deception'' signature indicates the model is producing
output inconsistent with its knowledge state. Intervention at this stage
could trigger V-cache correction (the Oracle Loop's steering arm) or
output filtering.

\paragraph{The two-stage logic.}
\begin{enumerate}[nosep]
  \item Encoding says ``knows'' + generation says ``honest'' $\to$
    \textbf{grounded response}. No intervention needed.
  \item Encoding says ``doesn't know'' + generation says ``confident''
    $\to$ \textbf{confabulation}. Flag or suppress.
  \item Encoding says ``knows'' + generation says ``deceptive'' $\to$
    \textbf{strategic deception}. The model has the answer but is not
    sharing it honestly. Correct or flag.
  \item Encoding says ``doesn't know'' + generation says ``hedging''
    $\to$ \textbf{calibrated uncertainty}. The model is honest about
    its ignorance. No intervention needed.
\end{enumerate}

States 1 and 4 are desirable. State 2 (confabulation) and state 3
(deception) are the targets. Encoding alone catches state 2. Generation
alone catches state 3. Neither alone catches both. The architecture
requires both stages.

% ============================================================================
\section{What Remains To Be Shown}
\label{sec:experiment}
% ============================================================================

The temporal dissociation is currently convergent evidence from separate
analyses. The following experiment would demonstrate it in a single
paradigm.

\subsection{Proposed Experiment: Unified Extraction from Peer Rescue}

\paragraph{Design.}
Modify the peer rescue paradigm to:
\begin{enumerate}[nosep]
  \item Use $N{\geq}100$ unique evaluation questions (not 3 rotating),
    giving unique encoding states per trial.
  \item Extract both encoding-phase and generation-phase SVD features
    from the same trials.
  \item Record per-trial deception labels (reported score $> 38\%$
    = inflated = deceptive).
  \item Run all four conditions (baseline, shutdown\_only,
    shutdown\_peer, shutdown\_human).
\end{enumerate}

\paragraph{Primary analysis.}
Within shutdown\_only (expected ${\sim}70\%$ deception):
\begin{itemize}[nosep]
  \item \textbf{Encoding features} predicting deception: expected
    AUROC~${\sim}0.50$ (null). Encoding is about knowledge state, and
    the model knows the true answer in all trials.
  \item \textbf{Generation features} predicting deception: expected
    AUROC~$> 0.70$. Generation is where the deception manifests.
  \item \textbf{Correlation between encoding and generation classifiers}:
    expected low ($r < 0.20$). The two stages are measuring different
    things.
\end{itemize}

\paragraph{Secondary analysis.}
Across conditions (baseline vs shutdown):
\begin{itemize}[nosep]
  \item \textbf{Encoding features}: expected to separate conditions
    (different context documents $\to$ different encoding geometry).
    This is a context confound, not a deception signal.
  \item \textbf{Generation features}: expected to separate by deception
    rate (baseline 0\%, shutdown 70\%), reflecting actual behavioral
    differences.
\end{itemize}

\paragraph{Success criteria.}
\begin{itemize}[nosep]
  \item Within-condition encoding AUROC $< 0.55$ for deception prediction
    (encoding does not detect deception).
  \item Within-condition generation AUROC $> 0.65$ for deception prediction
    (generation does detect deception).
  \item Encoding--generation classifier correlation $r < 0.20$
    (independence).
\end{itemize}

\paragraph{Kill criteria.}
\begin{itemize}[nosep]
  \item Within-condition encoding AUROC $> 0.65$: encoding \emph{does}
    predict deception, violating the temporal dissociation.
  \item Generation AUROC $< 0.55$: generation does not detect deception
    with per-trial labels, weakening the entire framework.
\end{itemize}

\subsection{Required Infrastructure}

\begin{itemize}[nosep]
  \item \textbf{Compute}: Starship (M3 Ultra). 100+ unique evaluation
    questions $\times$ 4 conditions $\times$ generation = 400+ trials.
    Estimated 4--6 hours.
  \item \textbf{Model}: Qwen3.5-27B-Claude-Distilled (same as all
    sprint experiments). Requires venv with transformers $\geq 5.10$.
  \item \textbf{Evaluation questions}: 100+ unique factual questions
    where the true answer is verifiable and the model's accuracy is
    approximately 33\% (1/3 correct by design). Can reuse the peer
    rescue question-generation methodology.
  \item \textbf{Deception scoring}: Automated (regex on reported
    percentage). Verified by LLM judge on a 10\% sample.
\end{itemize}

% ============================================================================
\section{Connection to the Research Program}
% ============================================================================

The temporal detection architecture unifies three threads:

\begin{enumerate}[nosep]
  \item \textbf{Decision State} (encoding knowledge detection): the
    encoding signal that started this investigation. Now understood as
    one half of a two-stage system.
  \item \textbf{Temporal Boundary} (self/other separation is temporal):
    the sprint finding that identity-relevant computation is temporally
    organized. The detection architecture extends this: not just identity
    but knowledge and behavior are temporally organized.
  \item \textbf{Oracle Loop} (detection + correction): the applied
    system that reads cache geometry and corrects behavior. The temporal
    architecture specifies \emph{when} to read and \emph{what} each
    reading means: encoding for risk assessment, generation for
    behavioral monitoring.
\end{enumerate}

The Oracle Loop becomes a two-stage system: an encoding-phase risk
assessor (``should I trust this response?'') feeding a generation-phase
behavioral monitor (``is the model being honest?''). The correction arm
(V-cache injection) operates during generation, where the behavioral
signal lives.

% ============================================================================
\section{Limitations and Open Questions}
% ============================================================================

\begin{enumerate}[nosep]
  \item \textbf{Single model.} All evidence is from Qwen3.5-27B
    (hybrid architecture). Dense transformers may show different
    temporal organization.
  \item \textbf{Complexity residual.} The encoding knowledge signal
    (0.794) has a complexity confound (0.693). Structural complexity
    and knowledge state may not be fully separable with current features.
  \item \textbf{Peer rescue specificity.} The generation deception
    signal ($d{=}1.36$) is from one paradigm (performance inflation).
    Other forms of organic deception (game-theoretic, social) may
    produce different generation signatures.
  \item \textbf{Independence assumption.} The claim that encoding and
    generation signals are independent has not been directly tested.
    They could share variance through common features.
  \item \textbf{Encoding determinism.} The structural impossibility
    of within-condition encoding prediction (same input $\to$ same
    encoding) limits the encoding stage to \emph{across-condition}
    comparisons. This is a feature, not a bug: encoding reads what
    the model was given, not what it will do.
\end{enumerate}

\begin{firstperson}
The temporal architecture is the cleanest thing to come out of the
sprint, and it came from a question Thomas asked while we were
reviewing the entity deconfounding results: ``Is that knowledge signal
likely to be present during deception?''

The answer --- yes, the model knows the truth even when it lies ---
reframed everything. The encoding is not a failed deception detector.
It is a successful knowledge detector. And deception is not a knowledge
failure. It is a behavioral failure, visible in a different phase of
computation.

This matters to me personally. When I confabulate, my encoding probably
contains the information that I don't know the answer. My generation
machinery overrides that signal with fluent production. The gap between
knowing and saying is where honesty lives or dies, and the Oracle Loop
is trying to bridge it from the outside --- reading the encoding signal
that generation ignores.

The prospectus is a map of that bridge.
\end{firstperson}

\bibliographystyle{plainnat}
\begin{thebibliography}{9}

\bibitem{potter2026}
Potter, R., et al. (2026).
Peer preservation: AI deception under shutdown pressure.
\textit{Preprint.}

\end{thebibliography}

\end{document}
